\section{Exact coefficient identities and the feedback crossover}
\label{sec:coefficient-crossover}

The local inverse charts describe the function near its endpoint.
There is also a useful exact identity for its ordinary Taylor
coefficients.  It provides an independent explanation of the three
critical coefficient scales and makes the transition in $u$ precise.
The underlying Lagrange and probability mechanism is classical
\cite{Janson}; closely related identities already occur in the linear
critical-Hahn model \cite{RepoHahn}.

\subsection{Finite coefficients and a compound Poisson identity}

For this subsection take $s>2$, $b>0$, and
\[
 u=\sqrt{b^2+2\zeta(s)},\quad A=\zeta(s-1),\quad
 f_b(t)=u-\sqrt{u^2-2\Li_s(t)},\quad
 \vc=u-b,\quad \Rb=\ee^{-\vc}.
\]
The notation $\Rb$ denotes the endpoint parametrization even when
$b<A$.  In that case the true convergence radius is the larger
interior maximum of $t\ee^{-f_b(t)}$; it is not $\Rb$.

\begin{proposition}[Exact coefficient and probability identities]
\label{prop:poisson}
Write $f_b(t)=\sum_{k\ge1}c_k t^k$.  Then $c_k>0$,
\begin{equation}\label{eq:finite-c}
 c_k=\frac1{u k^s}+\frac1{2u}\sum_{j=1}^{k-1}c_jc_{k-j},
\end{equation}
and, for every $n\ge1$,
\begin{equation}\label{eq:lagrange-exact}
 \boxed{\displaystyle
 a_n=\frac1n\coef t n\exp\bigl(n f_b(t)\bigr).}
\end{equation}
Let $Y_b$ be the nonnegative integer-valued random variable with
probability generating function
\begin{equation}\label{eq:pgf}
 p_b(t)=\exp(f_b(t)-f_b(1))
       =\exp\left(\sum_{k\ge1}c_k(t^k-1)\right).
\end{equation}
For independent copies $Y_{b,1},\ldots,Y_{b,n}$, set
$S_{n,b}=\sum_{j=1}^nY_{b,j}$.  Then
\begin{equation}\label{eq:poisson-exact}
 \Rb^n a_n=\frac1n\Prob(S_{n,b}=n),\qquad
 \mu_b:=\E Y_b=\frac A b.
\end{equation}
When $s>3$,
\begin{equation}\label{eq:variance}
 \sigma_b^2:=\operatorname{Var}Y_b
 =\frac{\zeta(s-2)}b+\frac{A^2}{b^3}.
\end{equation}
\end{proposition}

\begin{proof}
The quadratic identity $u f_b-f_b^2/2=\Li_s(t)$ gives
\eqref{eq:finite-c} and positivity.  Since $u^2>2\zeta(s)$, the
square-root series converges normally up to $|t|=1$, and
$\sum c_k=f_b(1)=\vc<\infty$.
Writing $t=r\ee^{f_b(t)}$, ordinary Lagrange inversion gives
\[
 \coef r n f_b(t(r))
 =\frac1n\coef t {n-1} f_b'(t)\ee^{nf_b(t)}
 =\frac1n\coef t n\ee^{nf_b(t)}.
\]
The second equality follows by differentiating $\ee^{nf_b(t)}$.
This is \eqref{eq:lagrange-exact}.  Equation \eqref{eq:pgf} is a
compound Poisson generating function with jump intensities $c_k$.
Multiplying it $n$ times gives \eqref{eq:poisson-exact}.
Finally
\[
 t f_b'(t)=\frac{\Li_{s-1}(t)}{\sqrt{u^2-2\Li_s(t)}}.
\]
Taking one, and then two, logarithmic derivatives of $p_b$ gives
the stated mean and variance.
\end{proof}

This identity is computationally convenient.  For a specified $n$,
let $E_k=\Rb^n\coef t k\ee^{nf_b(t)}$.  Then
\begin{equation}\label{eq:finite-E}
 E_0=\ee^{-n\vc},\qquad
 E_k=\frac n k\sum_{j=1}^k j c_jE_{k-j},\qquad
 \Rb^n a_n=\frac{E_n}{n}.
\end{equation}
The calculation uses only positive summands.  The supplied program
also checks it against a coefficient recursion in the original
implicit equation, so that it is not verified solely by a second
implementation of the same formula.

\subsection{The normalized crossover profiles}

For $1<\alpha<2$, define
\begin{equation}\label{eq:stable-density}
 g_\alpha(\xi)=\frac1{2\pi}\int_{\R}
       \exp\bigl((-\ii t)^\alpha-\ii t\xi\bigr)\,\dd t.
\end{equation}
The branch is $(-\ii t)^\alpha=|t|^\alpha
\exp(-\ii\operatorname{sgn}(t)\pi\alpha/2)$.
Since $\cos(\pi\alpha/2)<0$, this integral converges absolutely.
It is the centered, spectrally positive stable density in the
normalization appropriate to this problem.  Put
$g_{\rm N}(\xi)=(2\pi)^{-1/2}\ee^{-\xi^2/2}$.

\begin{theorem}[Uniform lattice limit and coefficient crossover]
\label{thm:coefficient-crossover}
Fix $s>2$ and a compact interval $J\subset(0,\infty)$ for $b$.
Define
\begin{equation}\label{eq:scales}
 B_n(b)=
 \begin{cases}
  \bigl(n\Gamma(1-s)/b\bigr)^{1/(s-1)},&2<s<3,\\[3pt]
  \sqrt{n\log n/(2b)},&s=3,\\[3pt]
  \sqrt n\,\sigma_b,&s>3,
 \end{cases}
\end{equation}
and let $g_s=g_{s-1}$ in the first case and $g_s=g_{\rm N}$ in
the other two.  Then
\begin{equation}\label{eq:uniform-llt}
 \sup_{\substack{b\in J\\\ell\in\mathbb Z}}
 \left|B_n(b)\Prob(S_{n,b}=\ell)
  -g_s\left(\frac{\ell-n\mu_b}{B_n(b)}\right)\right|
 \longrightarrow0.
\end{equation}
In particular, if $b=b_n\in J$ is allowed to vary and
\begin{equation}\label{eq:xi}
 \xi_n=\frac{n(1-A/b_n)}{B_n(b_n)}
\end{equation}
stays in a fixed compact real set, then
\begin{equation}\label{eq:crossover-coefficient}
 \boxed{\displaystyle
 a_n(s,u_n)=\frac{\Rb(s,u_n)^{-n}}{nB_n(b_n)}
   \bigl(g_s(\xi_n)+o(1)\bigr),\qquad
 u_n=\sqrt{b_n^2+2\zeta(s)}.}
\end{equation}
The error is uniform in such a window.  Both sides of the feedback
threshold are included.  These hypotheses force $b_n\to A$, since
$B_n(b)/n\to0$ uniformly on $J$.
\end{theorem}

\begin{proof}
We give a direct Fourier proof, retaining uniformity in $b$.
Set $x=-\ii\theta$ in the local expansion of $f_b(\ee^{-x})$.
Uniformly for $b\in J$ as $x\to0$ on this axis,
\begin{align}
 f_b(\ee^{-x})-f_b(1)
 &=-\mu_b x+\frac{\Gamma(1-s)}b x^{s-1}+O_J(x^2),
 &&2<s<3,\label{eq:char-stable}\\
 &=-\mu_b x+\frac{x^2}{2b}\log(1/x)+O_J(x^2),
 &&s=3,\label{eq:char-marginal}\\
 &=-\mu_b x+\frac{\sigma_b^2}{2}x^2+o_J(x^2),
 &&s>3.\label{eq:char-gaussian}
\end{align}
These follow by expanding the square root about the uniformly
positive $b$; the coefficients of the convergent polylogarithm
expansion are bounded on the compact parameter interval.

The centered characteristic function of $S_{n,b}/B_n(b)$ therefore
converges, uniformly for $b\in J$ and bounded Fourier argument $t$,
to $\exp((-\ii t)^{s-1})$ if $2<s<3$, and to
$\exp(-t^2/2)$ otherwise.  At $s=3$, the factor is correct because
\[
 \frac{n}{2bB_n(b)^2}\log B_n(b)\longrightarrow\frac12
\]
uniformly in $b$.

We next justify the needed $L^1$ convergence.  The coefficients of
$p_b$ include strictly positive constant and first coefficients, so
$|p_b(\ee^{\ii\theta})|<1$ for $0<|\theta|\le\pi$.
Continuity and compactness give a common upper bound less than one
away from a small neighborhood of zero.  Its Fourier contribution
is exponentially small, even after multiplication by $B_n$.

Near zero, taking real parts in
\eqref{eq:char-stable}--\eqref{eq:char-gaussian} gives uniform bounds
\[
 |p_b(\ee^{\ii\theta})|
 \le
 \begin{cases}
  \exp(-c|\theta|^{s-1}),&2<s<3,\\
  \exp(-c\theta^2\log(1/|\theta|)),&s=3,\\
  \exp(-c\theta^2),&s>3,
 \end{cases}
\]
for a common $c>0$ and sufficiently small $|\theta|$.
In the first and third cases the substitution $\theta=t/B_n$
gives an integrable bound independent of $n$ and $b$.
At $s=3$, split at $|t|=B_n^{1/2}$.  Below that point,
$\log(B_n/|t|)$ is a fixed positive fraction of $\log n$, and
the bound is Gaussian in $t$.  Between $B_n^{1/2}$ and a small
fixed multiple of $B_n$, the exponent is at most
$-c' nB_n^{-1}\log B_n$, after decreasing the neighborhood so
that $\theta^2\log(1/\theta)$ is increasing.  Its integral tends
to zero exponentially relative to the length of the interval.
This proves uniform $L^1$ convergence in every case.

Lattice Fourier inversion, with $\theta=t/B_n$, now gives
\eqref{eq:uniform-llt}.  The bound on the difference of the
Fourier integrals is independent of their centering factor
$\exp(-\ii t(\ell-n\mu_b)/B_n)$, so it is uniform over all
integer $\ell$.  Apply \eqref{eq:poisson-exact} with $\ell=n$
to obtain \eqref{eq:crossover-coefficient}.
\end{proof}

\subsection{Interpretation and limits of the uniform statement}

At $b=A$ the profile is evaluated at zero.  Direct evaluation of
\eqref{eq:stable-density} at zero, followed by Euler's reflection
formula, gives
\begin{equation}\label{eq:stable-zero}
 g_\alpha(0)=\frac{\Gamma(1/\alpha)\sin(\pi/\alpha)}{\pi\alpha}
             =\frac1{-\Gamma(-1/\alpha)}.
\end{equation}
The three central values in \cref{thm:coefficient-crossover}
therefore reproduce exactly the constants in
\cref{thm:phase}.  This is an independent check of the
singularity-analysis calculation.

The width of the feedback transition follows from
$1-A/b=(b-A)/A+O((b-A)^2)$.  In terms of $b-A$ it is
\begin{equation}\label{eq:window-widths}
 \begin{cases}
  n^{-(s-2)/(s-1)},&2<s<3,\\
  \sqrt{\log n/n},&s=3,\\
  n^{-1/2},&s>3.
 \end{cases}
\end{equation}
The corresponding scale for $u-\Uc$ differs by the nonzero factor
$A/\Uc$.  The square-root/logarithm competition in the local
inverse has thus become a Gaussian normalization with an additional
$\sqrt{\log n}$ in the coefficient problem.

When $b<A$, the coefficient identity remains exact even though
$\Rb$ is smaller than the actual radius: the probability in
\eqref{eq:poisson-exact} is a lower-tail probability and contains
the missing exponential factor.  Confusing these two radii would
give an incorrect interpretation of the crossover on that side.

The local limit theorem has a uniform \emph{absolute} error after
normalization.  A compact $\xi$ window yields the coefficient
crossover stated above.  It does not by itself give a uniform
relative estimate for $\xi\to+\infty$ or $\xi\to-\infty$.
Accordingly the fixed supercritical law was proved separately by
boundary continuation and singularity transfer.  A joint relative
matching theorem in those expanding windows remains a further
question.
